\chapter{Fever}

\section*{Definition and Overview}
Fever is a rise in body temperature above the normal range, usually defined as 38 degrees Celsius or higher, in response to infection or inflammation. It is not an illness itself but a sign that the body's immune system has been activated; the brain's thermostat is deliberately reset higher to help fight infection, and most fevers are brief and harmless. Fever is one of the most common reasons people, especially parents of young children, seek medical advice. The height of a fever is not a reliable guide to how serious the underlying illness is: how the person looks and behaves matters far more. This chapter explains what fever is, how to measure it, when to treat it, and when to seek urgent help.

\section*{Epidemiology}
Fever is extremely common. Most children under the age of 5 have several febrile illnesses a year, usually from viral infections that resolve on their own. Adults also develop fevers with viral and bacterial infections. While most fevers are benign, fever is the presenting feature of serious bacterial infections such as sepsis, meningitis, and pneumonia, and it is a leading reason for emergency department visits. Fevers from infectious causes are more common in winter, and the pattern of specific illnesses varies with age, season, and vaccination status. Most people who see a doctor with fever have a self-limiting viral illness.

\section*{Aetiology and Pathophysiology}
\begin{itemize}
    \item \textbf{Infections:} the most common cause; viruses (colds, influenza, gastroenteritis), bacteria (ear, chest, urine, and throat infections), and less commonly other organisms
    \item \textbf{Mechanism:} immune cells detect infection and release pyrogens such as interleukin-1, which signal the brain's hypothalamus to raise the temperature set point
    \item \textbf{Purpose:} the higher temperature slows some microbes and enhances the immune response
    \item \textbf{Inflammation:} autoimmune diseases, inflammatory conditions, and some cancers can cause fever without infection
    \item \textbf{Immunisation:} vaccines can cause short-lived fever as the immune system responds
    \item \textbf{Other causes:} heatstroke, some medicines, and hormonal conditions raise temperature through different mechanisms
    \item \textbf{Temperature regulation:} the body produces more heat and conserves it through shivering and reduced sweating until the set point is reached
    \item \textbf{Febrile seizures:} some young children convulse with rapid temperature rises, an alarming but usually benign event
\end{itemize}

\section*{Clinical Features}
\begin{itemize}
    \item A measured temperature of 38 degrees or higher, with normal temperature around 36--37.2 degrees
    \item Feeling hot, flushed, and sweaty, or shivering and cold during the rise phase
    \item Headache, muscle aches, fatigue, and loss of appetite
    \item In children: irritability, clinginess, reduced interest in play, and sometimes poor feeding
    \item The pattern matters: how long the fever has lasted, how high it goes, and whether it responds to paracetamol or ibuprofen
    \item The key feature for judging severity: the person's overall appearance, alertness, breathing, and ability to drink
    \item Accompanying symptoms point to the source: cough, runny nose, sore throat, ear pain, vomiting, diarrhoea, or urinary symptoms
\end{itemize}

\section*{Differential Diagnosis}
The cause of a fever is usually clear from the associated symptoms.
\begin{itemize}
    \item Viral infections: influenza, colds, gastroenteritis, and childhood exanthems
    \item Bacterial infections: ear infections, tonsillitis, pneumonia, urinary tract infection, and skin infections
    \item Serious infections: sepsis, meningitis, and kidney infection, which must be considered in unwell patients
    \item COVID-19 and other respiratory viruses
    \item Non-infectious causes: inflammatory diseases, heatstroke, some medicines, and rarely malignancy
    \item Post-immunisation fever
    \item In babies under 3 months, fever is treated seriously because the immune system is immature and serious infection is harder to detect
\end{itemize}

\section*{When to Seek Medical Care}
Most fevers can be managed at home, but medical review is needed when the person looks unwell, the fever lasts more than a few days, or there are worrying features. Babies under 3 months with a temperature of 38 degrees or more should be seen urgently by a doctor.

\textbf{Call 000 (or local emergency services) for:}
\begin{itemize}
    \item A baby under 3 months with a fever
    \item A child who is floppy, unresponsive, or difficult to wake
    \item Difficulty breathing, rapid breathing, or a blue tinge to the lips
    \item A non-blanching rash (one that does not fade when pressed), which can indicate meningococcal disease
    \item A seizure (fit)
    \item Signs of severe dehydration: no urine for 6--8 hours, sunken eyes, or a dry mouth
    \item Severe headache with neck stiffness, or confusion
    \item A fever in an adult with serious underlying illness, a weakened immune system, or who is very unwell
\end{itemize}

\section*{Diagnosis and Investigations}
\begin{itemize}
    \item \textbf{Measurement:} a digital thermometer under the arm, in the mouth (adults), or in the ear; in babies, a rectal thermometer is most accurate
    \item \textbf{History and examination:} looking for the source of infection in the ears, throat, chest, skin, and urine
    \item \textbf{Urine test:} in young children with unexplained fever, to detect urinary tract infection
    \item \textbf{Blood tests:} full blood count, inflammatory markers, and blood cultures in unwell or hospitalised patients
    \item \textbf{Imaging:} chest X-ray for suspected pneumonia
    \item \textbf{Lumbar puncture:} for suspected meningitis in unwell infants and patients with signs of meningism
    \item \textbf{In babies under 3 months:} a full sepsis work-up is standard because serious infection is common and hard to detect clinically
\end{itemize}

\section*{Management}
\begin{itemize}
    \item \textbf{Comfort measures:} dress lightly, keep the room comfortable, and encourage fluids; do not over-wrap or sponge with cold water
    \item \textbf{Paracetamol or ibuprofen:} used to make the person comfortable, not to chase a normal temperature; use one medicine at a time and follow the correct dose for age and weight
    \item \textbf{Do not use aspirin in children}, because of the risk of Reye's syndrome
    \item \textbf{Treat the cause:} antibiotics for bacterial infections such as ear, chest, or urinary infections; viral infections need no specific treatment
    \item \textbf{Watchful monitoring:} record the temperature, fluids, urine output, and the person's overall condition
    \item \textbf{Febrile seizures:} simple febrile seizures are managed with reassurance and first aid; emergency medicines are only for the small minority with prolonged seizures
    \item \textbf{When to review:} see a doctor if fever persists beyond 3 days, the person worsens, or new symptoms appear
\end{itemize}

\section*{Complications}
\begin{itemize}
    \item Dehydration from reduced intake and increased fluid loss, especially in young children
    \item Febrile seizures in susceptible young children
    \item Progression of the underlying infection to sepsis, pneumonia, or meningitis if not recognised
    \item Unnecessary antibiotics when fever is wrongly assumed to be bacterial
    \item Overdose from giving too much paracetamol or ibuprofen, or giving them too frequently
    \item Rarely, heat-related problems from over-treating or over-wrapping
\end{itemize}

\section*{Prognosis and Outcomes}
The vast majority of fevers settle within a few days as the underlying infection resolves, and the fever itself does no harm. In children, the height of the fever does not predict severity, and most febrile children recover fully without treatment of the fever itself. Serious bacterial infections, when present, respond well to prompt antibiotics, which is why the assessment focuses on identifying the unwell patient rather than the height of the temperature. With sensible monitoring and clear guidance on when to seek help, fever is managed safely at home in nearly all cases.

\section*{Prevention and Self-Care}
\begin{itemize}
    \item Keep a digital thermometer and age-appropriate paracetamol or ibuprofen at home, and know the correct doses
    \item Encourage fluids frequently; in babies, continue breast or formula feeds
    \item Dress the person lightly and keep the room cool but not cold
    \item Use paracetamol and ibuprofen only for comfort, at the correct dose and interval, and never together at the same time
    \item Wash hands and follow immunisation schedules to prevent common causes of fever
    \item Do not wrap a shivering child in extra blankets once the fever is up
    \item Seek advice early for babies under 3 months, people with chronic illness, or any fever lasting more than 3 days
\end{itemize}

\section*{Sources and Further Reading}
The following reputable sources provide further information for healthcare students and patients seeking reliable, current advice about fever.
\begin{itemize}
    \item Healthdirect Australia. \textit{Fever.} https://www.healthdirect.gov.au/fever
    \item The Royal Children's Hospital Melbourne. \textit{Fever in children fact sheet.} https://www.rch.org.au/
    \item Raising Children Network. \textit{Fever in babies and children.} https://raisingchildren.net.au/
    \item NHS. \textit{Fever in adults and children.} https://www.nhs.uk/conditions/fever-in-adults/
    \item Mayo Clinic. \textit{Fever.} https://www.mayoclinic.org/
    \item MedlinePlus. \textit{Fever.} https://medlineplus.gov/
\end{itemize}
